\section{Wigner--Weisskopf decay}\label{app:D.3}

This section derives the exponential decay of a discrete state coupled to a continuum, \cref{eq:2.4-WW}, the Lorentzian distribution of the final states, \cref{eq:2.4-lorentz}, and the replacement of an infinitesimal by a width in a second-order denominator, \cref{eq:2.4-width}. The original treatment is that of Weisskopf and Wigner~\cite{weisskopf1930berechnung}; we follow the resolvent method of Ref.~\cite[Chap.~III]{cohentannoudji1998atom}, and use the notation of \cref{sec:2.1,sec:2.4}.

\subsection*{Exact amplitude equation}

A discrete state $|i\rangle$ of energy $E_i$ is coupled by a constant perturbation $H'$ to a set of states $|f\rangle$ of energies $E_f$; the states $|f\rangle$ become a continuum with density $\mathcal{D}_f(E)$ as the volume grows. We assume $H'_{ii} = 0$ and $H'_{ff'} = 0$: the perturbation only connects $i$ with the continuum. A nonzero $H'_{ii}$ is a first-order shift and can be put into $E_i$; couplings within the continuum are left out of this model. With $\xi_f = E_f - E_i$, the amplitude equations \cref{eq:2.1-exact} are
\begin{equation}
  \ii\,\dot c_i(t) = \sum_f H'_{if}\,\ee^{-\ii\xi_ft}\,c_f(t), \qquad
  \ii\,\dot c_f(t) = H'_{fi}\,\ee^{\ii\xi_ft}\,c_i(t),
  \label{eq:D.3-amp}
\end{equation}
with $c_i(0) = 1$ and $c_f(0) = 0$. The second equation integrates to $c_f(t) = -\ii H'_{fi}\int_0^t\dd t'\,\ee^{\ii\xi_ft'}c_i(t')$. Inserting this solution into the first equation gives a closed equation for $c_i$,
\begin{equation}
  \begin{aligned}
  \dot c_i(t) &= -\int_0^t\dd t'\,K(t - t')\,c_i(t'), \\
  K(t - t') &= \sum_f|H'_{fi}|^2\,\ee^{-\ii\xi_f(t - t')} = \int\dd E\,\mathcal{D}_f(E)\,\overline{|H'_{fi}|^2}(E)\,\ee^{-\ii(E - E_i)(t - t')} .
  \end{aligned}
  \label{eq:D.3-kernel}
\end{equation}
No approximation has been made. The kernel $K$ is the Fourier transform of the product $\mathcal{D}_f\overline{|H'_{fi}|^2}$, which varies on the energy scale $E_{\mathrm{w}}$ of \cref{sec:2.1}. The kernel therefore decays on the time scale $1/E_{\mathrm{w}}$: the continuum has a memory of that length. For the flat continuum of \cref{we:2-continuum}, with $\mathcal{D}_f\overline{|H'_{fi}|^2} = \varGamma_i/2\pi$ on $|E - E_i| < E_{\mathrm{w}}/2$ and zero outside, $K(t - t') = (\varGamma_i/\pi)\sin[E_{\mathrm{w}}(t - t')/2]/(t - t')$, which starts at $\varGamma_iE_{\mathrm{w}}/2\pi$ and oscillates away within a few $1/E_{\mathrm{w}}$.

\subsection*{The Markov approximation}

At times short compared with $1/E_{\mathrm{w}}$, $c_i$ has barely changed and $\dot c_i \approx -t\,K(0)$, so $|c_i|^2 \approx 1 - K(0)t^2$: the decay starts quadratically, with the short-time behavior found in \cref{we:2-continuum}. At later times, the integral in \cref{eq:D.3-kernel} only receives contributions from $t - t' \lesssim 1/E_{\mathrm{w}}$. Suppose that $c_i$ changes little over that interval, as it does when its decay rate is small compared with $E_{\mathrm{w}}$. We can then replace $c_i(t')$ by $c_i(t)$ and extend the lower limit of the integral to $-\infty$:
\begin{equation*}
  \dot c_i(t) = -c_i(t)\int_{-\infty}^{t}\dd t'\,K(t - t') .
\end{equation*}
The remaining integral is a sum of terms $\int_{-\infty}^{t}\dd t'\,\ee^{-\ii\xi_f(t - t')}$, which converge with a factor $\ee^{-0^+(t - t')}$ and give $-\ii/(\xi_f - \ii0^+) = \ii/(E_i - E_f + \ii0^+)$. With the Sokhotski--Plemelj formula \cref{eq:2.4-SP},
\begin{equation}
  \begin{split}
  \int_{-\infty}^{t}\dd t'\,K(t - t') &= \ii\sum_f\frac{|H'_{fi}|^2}{E_i - E_f + \ii0^+}\\
  &= \ii\,\operatorname{PV}\sum_f\frac{|H'_{fi}|^2}{E_i - E_f} + \pi\sum_f|H'_{fi}|^2\,\delta(E_f - E_i)
  = \ii\Delta E_i + \frac{\varGamma_i}{2} ,
  \end{split}
  \label{eq:D.3-markov}
\end{equation}
where $\varGamma_i$ is the golden-rule rate \cref{eq:2.1-total} and $\Delta E_i$ the level shift of \cref{eq:2.4-WW}. The equation $\dot c_i = -(\varGamma_i/2 + \ii\Delta E_i)c_i$ integrates to \cref{eq:2.4-WW}. The Markov approximation requires $\varGamma_i \ll E_{\mathrm{w}}$, the condition under which the golden-rule window of \cref{eq:2.1-window} exists, and holds for $t \gg 1/E_{\mathrm{w}}$. The memory of the continuum is then short compared with the lifetime, and the probability to decay in the next instant is the same at every time; a constant decay probability per instant defines an exponential law.

\subsection*{The resolvent and the corrections}

The size of the corrections to the Markov approximation follows from an exact solution of \cref{eq:D.3-kernel} by Laplace transform. With $\tilde c_i(s) = \int_0^\infty\dd t\,\ee^{-st}c_i(t)$ for $\Real s > 0$, the convolution becomes a product, $s\tilde c_i - 1 = -\tilde K(s)\tilde c_i$, with $\tilde K(s) = \sum_f|H'_{fi}|^2/(s + \ii\xi_f)$. Define an energy variable $z = \ii s$, measured from $E_i$; for $\Real s > 0$ it lies in the upper half plane. Then
\begin{equation}
  \tilde c_i = \frac{\ii}{z - R(z)}, \qquad R(z) = \sum_f\frac{|H'_{fi}|^2}{z - \xi_f} ,
  \label{eq:D.3-resolvent}
\end{equation}
and $c_i(t) = (2\pi)^{-1}\int\dd z\,\ee^{-\ii zt}\,\ii/(z - R(z))$ along a line just above the real axis~\cite[Chap.~III]{cohentannoudji1998atom}. The function $R(z)$ is analytic off the real axis; just above the axis, by \cref{eq:2.4-SP}, $R(x + \ii0^+) = \operatorname{PV}\sum_f|H'_{fi}|^2/(x - \xi_f) - \ii\pi\sum_f|H'_{fi}|^2\delta(x - \xi_f)$, which at $x = 0$ is $\Delta E_i - \ii\varGamma_i/2$. When $R$ is continued below the axis through the continuum, $z - R(z)$ has a zero near $z_{\mathrm{p}} = \Delta E_i - \ii\varGamma_i/2$, as long as $R$ changes little over a range $\varGamma_i$ around $z = 0$. Closing the contour in the lower half plane picks up the pole of $\tilde c_i$ at this zero and gives
\begin{equation*}
  c_i(t) = r_{\mathrm{p}}\,\ee^{-\ii z_{\mathrm{p}}t} + c_i^{\mathrm{cut}}(t), \qquad r_{\mathrm{p}} = \frac{1}{1 - R'(z_{\mathrm{p}})} ,
\end{equation*}
plus the contribution $c_i^{\mathrm{cut}}$ of the branch points at the edges of the continuum. Since $R'$ is of order $\varGamma_i/E_{\mathrm{w}}$, the residue $r_{\mathrm{p}}$ of the pole differs from one by a relative amount of that order, and so does the decay rate, which becomes $\varGamma_i/(1 - R')$ when $R$ is expanded to first order about $z = 0$. For the flat continuum, $R(z) = (\varGamma_i/2\pi)\ln[(z + E_{\mathrm{w}}/2)/(z - E_{\mathrm{w}}/2)]$; just above $z = 0$ the function $R$ equals $-\ii\varGamma_i/2$, so the level shift vanishes, and $R'(0) = 2\varGamma_i/(\pi E_{\mathrm{w}})$. In \cref{we:2-continuum}, $\varGamma_i/E_{\mathrm{w}} = 0.01$ and $R'(0) = 0.0064$. The survival probability is then $r_{\mathrm{p}}^2\ee^{-\varGamma_it/(1 - R')}$, which at $t = 1/\varGamma_i$ exceeds $\ee^{-1}$ by $0.0023$, or \SI{0.4}{\percent} of the depletion. The exact solution for the flat continuum of \cref{we:2-continuum} deviates from the depletion $1 - \ee^{-\varGamma_it}$ by this amount once the short-time transient has died out. The branch-point contribution falls off as a power of $1/(E_{\mathrm{w}}t)$ with a prefactor of order $\varGamma_i/E_{\mathrm{w}}$. It overtakes the exponential only after many lifetimes, when the survival probability is already exponentially small, and it affects no rate in this book.

\subsection*{The line shape}

With \cref{eq:2.4-WW}, the amplitude of a final state is
\begin{equation*}
  c_f(t) = -\ii H'_{fi}\int_0^t\dd t'\,\ee^{\ii(\xi_f - \Delta E_i)t' - \varGamma_it'/2}
  \;\xrightarrow{\;t\to\infty\;}\; \frac{H'_{fi}}{E_f - E_i - \Delta E_i + \ii\varGamma_i/2} ,
\end{equation*}
whose square is \cref{eq:2.4-lorentz}. This Lorentzian is the natural line shape of the emitted radiation when the continuum consists of the photon modes~\cite{weisskopf1930berechnung}.

\subsection*{A decaying intermediate state}

Finally, let a state $|i\rangle$ be coupled weakly, by a constant $H_1$, to a state $|m\rangle$ that decays into a continuum through a second coupling $H_2$, with total rate $\varGamma_{\mathrm{w}}$ and shift $\Delta E_m$ from \cref{eq:D.3-markov} applied to $m$. Let $H_1$ act from the remote past, and keep it to first order, so that $c_i = 1$. The amplitude of $m$ then obeys
\begin{equation*}
  \dot c_m(t) = -\ii\,(H_1)_{mi}\,\ee^{\ii(E_m - E_i)t} - \Bigl(\frac{\varGamma_{\mathrm{w}}}{2} + \ii\Delta E_m\Bigr)c_m(t) ,
\end{equation*}
in which the decay term is the Markov form of the coupling to the continuum, valid under the conditions of the Markov approximation. The solution of this equation with $c_m(-\infty) = 0$ is
\begin{equation}
  c_m(t) = -\ii\,(H_1)_{mi}\int_{-\infty}^t\dd t_2\,\ee^{\ii(E_m - E_i)t_2}\,\ee^{-(\varGamma_{\mathrm{w}}/2 + \ii\Delta E_m)(t - t_2)}
  = \frac{(H_1)_{mi}\,\ee^{\ii(E_m - E_i)t}}{E_i - \tilde E_m + \ii\varGamma_{\mathrm{w}}/2} ,
  \label{eq:D.3-cm}
\end{equation}
with $\tilde E_m = E_m + \Delta E_m$. The free propagation $\ee^{-\ii E_m(t - t_2)}$ of \cref{eq:2.4-c2} has become a decaying one, and the integral converges without any switching factor. The amplitude of a final state $f$ reached from $m$ by $H_2$ is, at first order in $H_2$,
\begin{equation*}
  c_f(t) = -\ii\int_{-\infty}^t\dd t_1\,(H_2)_{fm}\,\ee^{\ii(E_f - E_m)t_1}\,c_m(t_1)
  = -\ii\,\frac{(H_2)_{fm}(H_1)_{mi}}{E_i - \tilde E_m + \ii\varGamma_{\mathrm{w}}/2}\int_{-\infty}^t\dd t_1\,\ee^{\ii(E_f - E_i)t_1} ,
\end{equation*}
which has the form of the second-order amplitude of \cref{sec:2.4} with the replacement \cref{eq:2.4-width}, and leads to \cref{eq:2.4-golden2} for the rate. The derivation needs a weak drive, $|(H_1)_{mi}| \ll |E_i - \tilde E_m + \ii\varGamma_{\mathrm{w}}/2|$, so that the population $|c_m|^2$ stays small and $i$ is not depleted; on resonance the condition reads $|(H_1)_{mi}| \ll \varGamma_{\mathrm{w}}/2$. The derivation also evaluates $\varGamma_{\mathrm{w}}$ and $\Delta E_m$ at the energy $E_m$ of $m$, whereas strictly the resolvent gives them at the energy $E_i$ at which $m$ is passed through. The values at $E_m$ and at $E_i$ agree when the decay continuum of $m$ varies little between $E_m$ and $E_i$, a condition that always holds near resonance, where the width matters.
